%% =====================================================================
%%  B3-T-09-01 -- what B3 contributes to "A Reputation You Built on Purpose"
%%  -------------------------------------------------------------------
%%  LAYER A: APPEND-ONLY. Written once, then left alone. Material found
%%  only here sits in the `unique` slot and is protected by rule R10.
%% =====================================================================
%% @kind: source
%% @id: B3-T-09-01
%% @book: B3
%% @topic: T-09-01
%% @locator: Law 5, pp. 58--67
%% @status: distilled
%% @unique: yes
%% @integrated: yes
%% @updated: 2026-09-19

\dossier{B3}{T-09-01}{\bookshort{B3}}{Law 5, pp. 58--67}

\begin{distilled}
  \item People judge on appearances --- clothes, gestures, words, actions --- because character is unfathomable and the visible is the only barometer available; one false slip can therefore be decisive.
  \item A reputation of your own creation protects you: it distracts probing eyes from what you are really like and gives you a measure of control over how the world judges you.
  \item Reputation is stated to work like magic: it can double your strength, and identical deeds read as brilliant or dreadful depending on the reputation of the doer.
  \item Attack form one: the showman case --- a rival's show was ruined by seeding doubts about his museum's stability; attendance collapsed and the show closed within weeks, while the saboteur built a lifelong name for audacity.
  \item Attack form two: once holding standing of his own, the same showman switched to gentle ridicule of rivals, which puts them on the defensive and draws attention to the mocker; outright slander is called too strong and liable to hurt the user.
  \item The law claims no reversal: neglecting reputation means letting others decide it for you, and even a reputation for insolence can be a deliberately chosen and valuable image.
\end{distilled}

\begin{unique}
  \item The denial dilemma stated as a mechanism: deny the rumour and a layer of suspicion remains (why defend so desperately?), ignore it and the doubt stands unrefuted --- the target pays under both responses, and defending can goad them into mistakes.
  \item The escalation rule: sow doubt from no standing at all; move to ridicule only from an established base; mockery signals security and double-serves as self-advertisement.
  \item The shield function: the published image absorbs the scrutiny that would otherwise reach the person --- reputation management as privacy.
\end{unique}
